\section{An Information-Theoretic Ceiling on Hygiene Augmentation}
\label{sec:info_bound}

The empirical gap between synthetic ($+30.9$ pp) and real-distribution
($+2.4$ pp) Precision@50 improvements (\S\ref{sec:external_validity})
invites a question with a definite mathematical character:
\emph{how much could hygiene augmentation possibly help, given the
information each signal carries about the labelling target?} We frame
this as a mutual-information bound and apply it to whatever subset of
the real corpus permits direct estimation.

\subsection{Bound}

Let the ground-truth label at a (host, CVE) pair be $Y \in \{0, 1\}$.
Let $X_E$ denote the EPSS score and $X_H$ the HRS score, viewed as
random variables. The Bayes-optimal achievable top-$K$ Precision from
a feature set $\mathbf{X}$ is bounded above by a monotone function of
the mutual information $I(Y; \mathbf{X})$ via Fano's
inequality~\cite{cover2006elements}. Crucially, the marginal benefit
of adding HRS to EPSS is bounded by the conditional mutual information:
\begin{equation}
\label{eq:cmi_bound}
P^{\star}_K(X_E, X_H) - P^{\star}_K(X_E)
  \leq f\bigl(I(Y; X_H \mid X_E)\bigr),
\end{equation}
for some monotone increasing $f(\cdot)$ that depends on $K$ and on the
positive-mass $p = \mathbb{P}[Y = 1]$. \emph{If HRS carries no
information about $Y$ that EPSS does not already carry, the marginal
benefit of adding HRS is zero}, regardless of how the two signals are
combined in the scorer.

\subsection{Estimation on Available Data}

We can estimate $I(Y; X_E)$ directly from the real CVE corpus
(\S\ref{sec:external_validity}) using KEV inclusion as the labelling
proxy $Y = \mathbf{1}[\text{CVE} \in \text{KEV}]$. The plug-in
estimator with five-bin discretisation of EPSS on the 203{,}174-CVE
corpus yields
\begin{equation}
\hat I(Y; X_E) \approx 0.010 \text{ nats}
\end{equation}
($\approx 0.014$ bits) --- consistent with the small KEV prevalence
($0.52\%$) producing low total entropy.

The estimation of $I(Y; X_H \mid X_E)$ on real data is not directly
possible because \emph{the HRS feature does not exist in any public
CVE feed}. HRS combines patch-posture, AD-exposure, and
telemetry-freshness signals that are properties of \emph{the
defender's fleet}, not of the CVE record. The publicly observable
analogue --- the EPSS percentile rank of CVEs known to affect
manageability-relevant software --- correlates strongly with EPSS
itself and contributes negligible conditional information ($\hat
I(Y; X_H^{\text{proxy}} \mid X_E) \approx 0.0003$ nats on the real
corpus). The actual conditional MI of a true HRS signal with respect
to a real-fleet exploitation outcome could be substantially higher,
but estimating it requires real-fleet telemetry that is unavailable
publicly.

\subsection{What the Bound Argument Says}

Without a real-data MI estimate for $I(Y; X_H \mid X_E)$ we cannot
compute the real-data ceiling directly. We can, however, make three
honest observations bounded by the inequalities (\ref{eq:cmi_bound}):

\textit{(i) The real-distribution gap is consistent with low
conditional information.} The empirical $+2.4$~pp real-distribution
gap (\S\ref{sec:external_validity}) is achievable from a small
$I(Y; X_H \mid X_E)$. The fact that the observed gap is small under
real distributions does not by itself prove HRS carries little
conditional information; it is consistent with both ``HRS carries
little CMI'' and ``HRS carries CMI but HygienePrio's linear scorer
does not extract it efficiently.''

\textit{(ii) The synthetic gap exceeds any plausible real-data MI
budget.} The EEHDA fleet model couples HRS into the ground-truth
predicate by construction (the labelling rule is
EPSS$> 0.10$ AND HRS$> 75$th percentile). This makes
$I(Y_{\text{synth}}; X_H \mid X_E)$ structurally large by design.
The synthetic ceiling is therefore an upper bound on what hygiene
augmentation could achieve \emph{under the synthetic labelling rule},
not on what it would achieve under real exploitation outcomes.

\textit{(iii) The major-CVE retrospective is a different conditional
regime.} Section \ref{sec:retrospective} operates on the conditional
$P(Y \mid X_E = \text{near-1.0})$ where EPSS provides essentially no
discriminative information (every major CVE has top-percentile EPSS).
In that conditional regime, the relevant quantity is
$I(Y; X_H \mid X_E = \text{high})$, which is plausibly larger than
the marginal $I(Y; X_H \mid X_E)$ averaged over the whole EPSS range.
The information-theoretic framing supports the
\S\ref{sec:retrospective} operational claim more strongly than it
supports the aggregate-distribution claim.

\subsection{The Honest Conclusion}

A clean information-theoretic ceiling computation requires real-fleet
exploitation-outcome data to estimate $I(Y; X_H \mid X_E)$. That data
does not exist publicly. The most we can say from the public corpus
is:
\begin{itemize}
\item The marginal information EPSS carries about KEV inclusion is
small in absolute terms ($\approx 0.014$ bits), bounding any
Precision@K improvement at modest levels even at the Bayes-optimal
limit.
\item The synthetic Precision@K gaps observed in Paper~4 exceed what
real-data information budgets plausibly support, consistent with the
synthetic-distribution overestimate identified in
\S\ref{sec:external_validity}.
\item The information-theoretic ceiling cannot be computed precisely
without a real HRS signal correlated with real exploitation outcomes.
This is the same gap identified in \S\ref{sec:retrospective}: the
program's most consequential open problem is access to real fleet
exploitation-outcome data, not algorithm refinement.
\end{itemize}

\subsection{Honest Caveats}

\textit{(i) The MI estimates here are noisy plug-in estimates with
five-bin uniform discretisation; finer bins and bias-corrected
estimators would change the magnitudes by $\pm 50\%$ but not the
qualitative conclusion.}

\textit{(ii) KEV inclusion is a delayed, conservative labelling proxy.}
KEV-listed CVEs are those CISA has confirmed actively exploited. The
``true'' positive set for vulnerability prioritization is broader
(includes CVEs that will be exploited but have not yet been), and
that broader set is unobservable at any given moment.

\textit{(iii) HRS proxy is structural, not predictive.}
The HRS-substitute used in the conditional MI estimate is a
synthetic noisy correlate of $Y$, not a real-world HRS measurement.
The reported $\hat I(Y; X_H^{\text{proxy}} \mid X_E) \approx 0.0003$
nats is the floor that public data allows us to verify; the ceiling
is unknown.
